\ifdefined\mainfile\else
\documentclass{article}
\usepackage{graphicx}
\usepackage{float}
\usepackage{amsmath}
\usepackage{amssymb}
\usepackage[a4paper, left=2.5cm, right=2.5cm, top=2.5cm, bottom=2.5cm]{geometry}
\usepackage[colorlinks=true,linkcolor=black,citecolor=blue,urlcolor=blue]{hyperref}
\begin{document}
\fi

\section{Reflection}
\label{sec:reflection}

The five biggest problems we ran into during the project, and how we resolved each:

\begin{itemize}
    \item \textbf{Pin conflicts on the ESP8266.} The ESP only has a handful of "safe" GPIOs and several of those have boot-time roles. Adding the OLED forced us to move a stepper coil pin (IN4 from D1 to D0), and a temporary MOSFET we put on the stepper power line turned out to be wired to TX, which silently killed serial logging until we tracked it down. We drew up a full pinout table on a whiteboard before adding the OLED, freed the I\textsuperscript{2}C bus by moving the coil, and replaced the MOSFET with a software trick -- pulling all four coil pins LOW after each move, which gives the same idle-current saving with one fewer component on the board.
    \item \textbf{Noisy LM35 readings.} Early temperature readings were both biased and noisy: the bias came from a miscalibrated ADC reference, and the noise from sampling at full rate without smoothing. Calibrating the reference to 2.667\,V and adding an IIR filter (10-sample running mean blended 30/70 with the previous value) brought the steady-state error to within $\pm$1\,\textdegree C.
    \item \textbf{Unreliable closure.} A single motor pulling one string caused the bag to twist on itself instead of sealing. Switching to two motors -- one on each side, each with its own string -- gave reliable closure across $>$20 attempts.
    \item \textbf{OFFLINE detection.} At first we checked the HTTP status code to decide whether the bin was online, but ThingSpeak still returns 200 OK when the ESP is unplugged because it keeps serving the last cached values. We changed it to look at the timestamp of the newest sensor reading instead: if that reading is more than 90 seconds old, the dashboard shows OFFLINE.
\end{itemize}
\noindent
Looking back, the thing that surprised us most was how much a free third-party service ended up dictating the design. ThingSpeak's free tier only allows one write every 15 seconds, so we set the push interval to 16 seconds, and that is the main reason a bag-count change made from the dashboard can take up to about a minute to reach the bin. We also underestimated the mechanical work: the string-closing mechanism was redesigned several times before it closed the bag reliably, while the firmware was comparatively quick to get working. What helped us the most day to day was keeping the debugging simple. The OFFLINE indicator on the dashboard and the serial log made it easy to see what the bin was doing whenever something broke.
\\ \\
A complete who-did-what breakdown and the percentage distribution of effort is given in the appendix.

\ifdefined\mainfile\else
\end{document}
\fi
